% Source: https://artofproblemsolving.com/community/c3h2351631_3_circles_and_a_triangle
% Stable identity: aops-topic-2351631-q01
Put $C=(0,0)$, $A=(3,0)$ and $B=(0,4)$. The blue circle has centre $\left(\dfrac32,0\right)$ and radius $\dfrac32$, the green circle centre $(0,2)$ and radius $2$, and the pink circle centre $\left(\dfrac32,2\right)$ and radius $\dfrac52$. Let the orange circle have centre $(x,y)$ and radius $r$. External tangency to the blue and green circles and internal tangency to the pink circle mean
\begin{align*}
\left(x-\dfrac32\right)^2+y^2&=\left(\dfrac32+r\right)^2,\\
x^2+(y-2)^2&=(2+r)^2,\\
\left(x-\dfrac32\right)^2+(y-2)^2&=\left(\dfrac52-r\right)^2.
\end{align*}
Subtracting the third equation from the first: $4y-4=\left(\dfrac94+3r\right)-\left(\dfrac{25}4-5r\right)=8r-4$, so $y=2r$. Subtracting the third from the second: $3x-\dfrac94=\left(4+4r\right)-\left(\dfrac{25}4-5r\right)=9r-\dfrac94$, so $x=3r$. Substituting into the third equation,
\[\left(3r-\dfrac32\right)^2+(2r-2)^2=\left(\dfrac52-r\right)^2\iff13r^2-17r+\dfrac{25}4=r^2-5r+\dfrac{25}4\iff12r(r-1)=0.\]
Since $r>0$, $r=1$, and the centre is $O=(3,2)$.

\begin{asy}
usepackage("amsmath");
size(6cm);
defaultpen(fontsize(10pt));
pair C=(0,0), A=(3,0), B=(0,4);
draw(circle((1.5,2),2.5),magenta+linewidth(0.9));
draw(circle((1.5,0),1.5),blue+linewidth(0.9));
draw(circle((0,2),2),heavygreen+linewidth(0.9));
draw(circle((3,2),1),orange+linewidth(1.1));
draw(A--B--C--cycle,linewidth(1));
draw((3,2)--(4,2),orange); draw((1.5,0)--(3,2),dashed+gray(0.4)); draw((0,2)--(3,2),dashed+gray(0.4));
dot(C); dot(A); dot(B); dot((3,2),orange); dot((1.5,0),blue); dot((0,2),heavygreen); dot((1.5,2),magenta);
label("$C$",C,SW); label("$A$",A,SE); label("$B$",B,NW); label("$O$",(3,2),N,orange);
label("$1$",(3.5,2),N,orange+fontsize(8pt));
label("$3$",(0.8,0),S,fontsize(8pt)); label("$4$",(0,3.1),W,fontsize(8pt));
\end{asy}

Check: $O$ is at distance $\dfrac52=\dfrac32+1$ from the blue centre, $3=2+1$ from the green centre and $\dfrac32=\dfrac52-1$ from the pink centre; it lies beyond the hypotenuse, in the region of the pink disc not covered by the other two discs. The radius is $\boxed{1}$.
